% Bagian: first-order-logic
% Bab: syntax-and-semantics
% Subbagian: formation-sequences

\documentclass[../../../include/open-logic-section]{subfiles}

\begin{document}

\olfileid[id]{fol}{syn}{fseq}

\olsection{Barisan Pembentukan}

Mendefinisikan !!{formula}s melalui suatu definisi induktif, beserta teknik
pelengkap berupa pembuktian sifat-sifat !!{formula}s melalui induksi,
merupakan pendekatan yang anggun dan efisien. Namun, pendekatan dari bawah
ke atas dan langkah demi langkah terhadap pembangunan !!{formula}s juga
dapat berguna. Di sini kita melakukannya dengan menggunakan gagasan
\emph{barisan pembentukan}.
%
Untuk menunjukkan bagaimana suku dan !!{formula}s dapat diperkenalkan dengan
cara ini tanpa perlu merujuk kepada definisi induktifnya, mula-mula kita
memperkenalkan gagasan untai simbol sembarang yang diambil dari suatu
bahasa~$\Lang L$.

\begin{defn}[Untai]
\ollabel{defn:string}
Andaikan $\Lang L$ merupakan suatu bahasa orde pertama. Suatu
\emph{untai-$\Lang L$} adalah barisan berhingga simbol-simbol dari~$\Lang L$.
Jika bahasa~$\Lang L$ jelas ditetapkan oleh konteks, kita akan sering
menyebut untai-$\Lang L$ cukup sebagai \emph{untai}.
\end{defn}

\begin{ex}
Untuk setiap bahasa orde pertama $\Lang L$, semua
!!{formula}s-$\Lang L$ merupakan untai-$\Lang L$, tetapi tidak sebaliknya.
Sebagai contoh, 
\[)(\Obj v_0\lif\lexists\]
merupakan untai-$\Lang L$, tetapi bukan !!{formula}-$\Lang L$.
\end{ex}

\begin{defn}[Barisan pembentukan untuk suku]
\ollabel{defn:fseq-trm}
Suatu barisan berhingga untai-$\Lang L$, yaitu
$\tuple{t_0,\dotsc,t_n}$, merupakan \emph{barisan pembentukan} untuk suatu
suku $t$ jika $t \ident t_n$ dan, untuk setiap $i \leq n$, $t_i$ merupakan
!!a{variable} atau !!a{constant}, atau $\Lang L$ memuat suatu !!{function}
$k$-tempat~$f$ dan terdapat $m_1,\dotsc,m_k < i$ sedemikian sehingga
$t_i \ident f(t_{m_1},\dotsc,t_{m_k})$. Jika perlu dibedakan, barisan
pembentukan untuk suku akan kita sebut \emph{barisan pembentukan suku}.
\end{defn}

\begin{ex}
Barisan
\[
    \tuple{\Obj c_0, \Obj v_0, \Atom{\Obj f^2_0}{\Obj c_0, \Obj v_0}, \Atom{\Obj f^1_0}{\Atom{\Obj f^2_0}{\Obj c_0, \Obj v_0}}}
\]
merupakan barisan pembentukan untuk suku $\Atom{\Obj f^1_0}{\Atom{\Obj
f^2_0}{\Obj c_0, \Obj v_0}}$, demikian pula
\[
    \tuple{\Obj v_0, \Obj c_0, \Atom{\Obj f^2_0}{\Obj c_0, \Obj v_0}, \Atom{\Obj f^1_0}{\Atom{\Obj f^2_0}{\Obj c_0, \Obj v_0}}}.
\]
\end{ex}

\begin{defn}[Barisan pembentukan untuk formula]
\ollabel{defn:fseq-frm}
Suatu barisan berhingga untai-$\Lang L$, yaitu
$\tuple{!A_0,\dotsc,!A_n}$, merupakan \emph{barisan pembentukan} untuk~$!A$
jika $!A \ident !A_n$ dan, untuk setiap $i \leq n$, $!A_i$ merupakan
!!{formula} atomik atau terdapat $j,k < i$ serta !!a{variable}~$x$
sedemikian sehingga salah satu dari hal berikut berlaku:
\begin{enumerate}
    \tagitem{prvNot}{$!A_i \ident \lnot !A_j$.}{}%
    \tagitem{prvAnd}{$!A_i \ident (!A_j \land !A_k)$.}{}%
    \tagitem{prvOr}{$!A_i \ident (!A_j \lor !A_k)$.}{}%
    \tagitem{prvIf}{$!A_i \ident (!A_j \lif !A_k)$.}{}%
    \tagitem{prvIff}{$!A_i \ident (!A_j \liff !A_k)$.}{}%
    \tagitem{prvAll}{$!A_i \ident \lforall[x][!A_j]$.}{}%
    \tagitem{prvEx}{$!A_i \ident \lexists[x][!A_j]$.}{}%
\end{enumerate}
Jika perlu dibedakan, barisan pembentukan untuk formula akan kita sebut
\emph{barisan pembentukan formula}.
\end{defn}

\begin{ex}
\[
\tuple{
    \Atom{\Obj A^1_0}{\Obj v_0},
    \Atom{\Obj A^1_1}{\Obj c_1},
    (\Atom{\Obj A^1_1}{\Obj c_1} \land \Atom{\Obj A^1_0}{\Obj v_0}),
    \lexists[\Obj v_0][(\Atom{\Obj A^1_1}{\Obj c_1} \land \Atom{\Obj A^1_0}{\Obj v_0})]
}
\]
merupakan barisan pembentukan untuk $\lexists[\Obj v_0][(\Atom{\Obj A^1_1}{\Obj
c_1} \land \Atom{\Obj A^1_0}{\Obj v_0})]$, demikian pula
\begin{multline*}
\tuple{
    \Atom{\Obj A^1_0}{\Obj v_0},
    \Atom{\Obj A^1_1}{\Obj c_1},
    (\Atom{\Obj A^1_1}{\Obj c_1} \land \Atom{\Obj A^1_0}{\Obj v_0}),
    \Atom{\Obj A^1_1}{\Obj c_1},\\
    \lforall[\Obj v_1][\Atom{\Obj A^1_0}{\Obj v_0}],
    \lexists[\Obj v_0][(\Atom{\Obj A^1_1}{\Obj c_1} \land \Atom{\Obj A^1_0}{\Obj v_0})]
}.
\end{multline*}
%
Sebagaimana tampak dari contoh kedua, barisan pembentukan dapat memuat
``sampah'': !!{formula}s yang mubazir atau tidak turut membangun hasil
akhirnya.
\end{ex}

\begin{prop}\ollabel{prop:formed}
Setiap !!{formula}~$!A$ dalam~$\Frm[L]$ memiliki barisan pembentukan.
\end{prop}

\begin{proof}
Andaikan $!A$ atomik. Maka barisan $\tuple{!A}$ merupakan barisan
pembentukan untuk~$!A$.
%
Sekarang andaikan $!B$ dan~$!C$ masing-masing memiliki barisan pembentukan
$\tuple{!B_0,\dotsc,!B_n}$ dan $\tuple{!C_0,\dotsc,!C_m}$.
%
\begin{enumerate}
  \tagitem{prvNot}{Jika $!A \ident \lnot !B$,
      maka $\tuple{!B_0,\dotsc,!B_n,\lnot !B_n}$
      merupakan barisan pembentukan untuk~$!A$.}{}
  \tagitem{prvAnd}{Jika $!A \ident (!B \land !C)$,
      maka $\tuple{!B_0,\dotsc,!B_n,!C_0,\dotsc,!C_m,(!B_n \land !C_m)}$
      merupakan barisan pembentukan untuk~$!A$.}{}
  \tagitem{prvOr}{Jika $!A \ident (!B \lor !C)$,
      maka $\tuple{!B_0,\dotsc,!B_n,!C_0,\dotsc,!C_m,(!B_n \lor !C_m)}$
      merupakan barisan pembentukan untuk~$!A$.}{}
  \tagitem{prvIf}{Jika $!A \ident (!B \lif !C)$,
      maka $\tuple{!B_0,\dotsc,!B_n,!C_0,\dotsc,!C_m,(!B_n \lif !C_m)}$
      merupakan barisan pembentukan untuk~$!A$.}{}
  \tagitem{prvIff}{Jika $!A \ident (!B \liff !C)$,
      maka $\tuple{!B_0,\dotsc,!B_n,!C_0,\dotsc,!C_m,(!B_n \liff !C_m)}$
      merupakan barisan pembentukan untuk~$!A$.}{}
  \tagitem{prvAll}{Jika $!A \ident \lforall[x][!B]$,
      maka $\tuple{!B_0,\dotsc,!B_n,\lforall[x][!B_n]}$
      merupakan barisan pembentukan untuk~$!A$.}{}
  \tagitem{prvEx}{Jika $!A \ident \lexists[x][!B]$,
      maka $\tuple{!B_0,\dotsc,!B_n,\lexists[x][!B_n]}$
      merupakan barisan pembentukan untuk~$!A$.}{}
\end{enumerate}
Menurut prinsip induksi pada !!{formula}s, setiap !!{formula} memiliki
barisan pembentukan.
\end{proof}

Kita juga dapat membuktikan arah sebaliknya. Hal ini penting karena
menunjukkan bahwa kedua cara kita mendefinisikan formula adalah ekuivalen:
keduanya memberikan hasil yang sama. Hal ini juga berarti bahwa kita dapat
membuktikan teorema tentang formula dengan menggunakan induksi biasa pada
panjang barisan pembentukan.

\begin{lem}
\ollabel{lem:fseq-init}
Andaikan $\tuple{!A_0,\dotsc,!A_n}$ merupakan barisan pembentukan
untuk~$!A_n$, dan $k \leq n$. Maka $\tuple{!A_0,\dotsc,!A_k}$ merupakan
barisan pembentukan untuk~$!A_k$.
\end{lem}

\begin{proof}
Latihan.
\end{proof}

\begin{prob}
Buktikan \olref[fol][syn][fseq]{lem:fseq-init}.
\end{prob}

\begin{thm}
\ollabel{thm:fseq-frm-equiv}
$\Frm[L]$ merupakan himpunan semua untai-$\Lang L$~$!A$ sedemikian sehingga
terdapat suatu barisan pembentukan formula untuk~$!A$.
\end{thm}

\begin{proof}
Misalkan $F$ ialah himpunan semua untai simbol dalam bahasa~$\Lang L$ yang
memiliki barisan pembentukan. Kita telah melihat dalam
\olref[fol][syn][fseq]{prop:formed} bahwa $\Frm[L] \subseteq F$, sehingga
sekarang kita membuktikan arah sebaliknya.

Andaikan $!A$ memiliki barisan pembentukan $\tuple{!A_0,\dotsc,!A_n}$. Kita
membuktikan bahwa $!A \in \Frm[L]$ melalui induksi kuat pada~$n$, yaitu
indeks terakhir barisan tersebut. Hipotesis induksi kita menyatakan bahwa
setiap untai simbol dengan barisan pembentukan berindeks terakhir $m < n$
berada dalam~$\Frm[L]$. Menurut definisi barisan pembentukan, identitas
$!A \ident !A_n$ berlaku dengan ruas kanannya atomik, atau harus terdapat
$j,k < n$ sedemikian sehingga salah satu dari hal berikut berlaku:
\begin{enumerate}
    \tagitem{prvNot}{$!A \ident \lnot !A_j$.}{}%
    \tagitem{prvAnd}{$!A \ident (!A_j \land !A_k)$.}{}%
    \tagitem{prvOr}{$!A \ident (!A_j \lor !A_k)$.}{}%
    \tagitem{prvIf}{$!A \ident (!A_j \lif !A_k)$.}{}%
    \tagitem{prvIff}{$!A \ident (!A_j \liff !A_k)$.}{}%
    \tagitem{prvAll}{$!A \ident \lforall[x][!A_j]$.}{}%
    \tagitem{prvEx}{$!A \ident \lexists[x][!A_j]$.}{}%
\end{enumerate}
Sekarang kita bernalar menurut kasus. Jika $!A$ atomik, maka
$!A_n \in \Frm[L]$. Sebaliknya, andaikan $!A \ident
(!A_j \land !A_k)$. Menurut \olref[fol][syn][fseq]{lem:fseq-init},
$\tuple{!A_0,\dotsc,!A_j}$ dan $\tuple{!A_0,\dotsc,!A_k}$ masing-masing
merupakan barisan pembentukan untuk $!A_j$ dan~$!A_k$. Karena keduanya
merupakan subbarisan awal sejati dari barisan pembentukan untuk~$!A$, indeks
terakhir keduanya lebih kecil daripada~$n$. Oleh karena itu, menurut
hipotesis induksi, $!A_j$ dan~$!A_k$ berada dalam~$\Frm[L]$, sehingga
menurut definisi !!a{formula}, $(!A_j \land !A_k)$ juga berada di dalamnya.
Kasus-kasus lain mengikuti penalaran yang sejajar.
\end{proof}

Barisan pembentukan untuk suku memiliki sifat-sifat yang serupa dengan
barisan pembentukan untuk !!{formula}s.

\begin{prop}
\ollabel{prop:fseq-trm-equiv}
$\Trm[L]$ merupakan himpunan semua untai-$\Lang L$~$t$ sedemikian sehingga
terdapat suatu barisan pembentukan suku untuk~$t$.
\end{prop}

\begin{proof}
Latihan.
\end{proof}

\begin{prob}
Buktikan \olref[fol][syn][fseq]{prop:fseq-trm-equiv}.
Petunjuk: gunakan strategi yang serupa dengan strategi dalam bukti
\olref[fol][syn][fseq]{thm:fseq-frm-equiv}.
\end{prob}

Ada dua jenis ``sampah'' yang dapat muncul dalam barisan pembentukan: unsur
yang berulang dan unsur yang tidak relevan bagi pembangunan formula atau
suku. Kita dapat menghilangkan keduanya dengan meninjau barisan pembentukan
minimal.

\begin{defn}[Barisan pembentukan minimal]
\ollabel{defn:minimal-fseq}
Suatu barisan pembentukan $\tuple{!A_0, \ldots, !A_n}$ untuk suatu
formula~$!A$ merupakan \emph{barisan pembentukan minimal} untuk~$!A$ jika,
untuk setiap barisan pembentukan lain~$s$ untuk~$!A$, panjang~$s$ lebih
besar daripada atau sama dengan~$n+1$.

Demikian pula, suatu barisan pembentukan $\tuple{t_0, \ldots, t_n}$ untuk
suatu suku~$t$ merupakan \emph{barisan pembentukan minimal} untuk~$t$ jika,
untuk setiap barisan pembentukan lain~$s$ untuk~$t$, panjang~$s$ lebih besar
daripada atau sama dengan~$n+1$.
\end{defn}

Perhatikan bahwa suatu formula atau suku dapat memiliki lebih dari satu
barisan pembentukan minimal, tetapi barisan-barisan tersebut akan memuat
tepat untai-untai yang sama.

\begin{prop}
\ollabel{prop:subformula-equivs}
Pernyataan-pernyataan berikut ekuivalen:
\begin{enumerate}
    \item $!B$ merupakan !!{subformula} dari~$!A$.
    \item $!B$ muncul dalam setiap barisan pembentukan untuk~$!A$.
    \item $!B$ muncul dalam suatu barisan pembentukan minimal untuk~$!A$.
\end{enumerate}
\end{prop}

\begin{proof}
Latihan.
\end{proof}

\begin{prob}
Buktikan \olref[fol][syn][fseq]{prop:subformula-equivs}.
\end{prob}

\begin{history}
Barisan pembentukan diperkenalkan oleh Raymond Smullyan dalam buku teksnya,
\emph{First-Order Logic} \citep{Smullyan1968}. Sifat-sifat tambahan barisan
pembentukan ditetapkan oleh \citet{Zuckerman1973}.
\end{history}

\end{document}
